\section{Engineering evidence and its availability}\label{sec:positive_results}

\subsection{Setting and research questions}

The existing EngiBench Beams2D v1 workflow \citep{ref23} uses a symmetric MBB beam with $100\times50$ bilinear elements, SIMP exponent 3, stiffness limits $10^{-9}$ and 1, and Poisson ratio 0.3. Nineteen effective logits and an anchor define a $4\times5$ coarse density field interpolated to the mesh; a logistic shift fixes mean density at 0.35. Smaller compliance is preferable. This is one deterministic structural-design family.

Both ten-group datasets retain 60 metadata states: 17 alias-only states and 43 nontrivial states containing 114 first-action branches, with eligible pools of 1024 designs. All have already been observed in development. The former holdout name does not create an unseen validation set for this method.

The existing evidence addresses three questions: whether finite-menu selection can meet the declared numerical tolerance; whether local stability reduces complete projection under the same objective and guarantee; and what the already paid two-evaluation engineering outcomes show. The relevant evidence is, respectively, checked enclosures, the direct same-guarantee control, and common-continuation paid records. This integration adds no experimental run.

\subsection{Availability before numerical interpretation}

Private-branch commit \texttt{7b36e70d8fb4d81d47bdd00a579453645c6828d6} preserves the historical positive-response aggregate. It reports completed execution and checking. Its 13 executed source files, per-state certificates and posterior enclosures, frozen decisions and raw timing records were not uploaded before a storage failure. A subsequent read-only check found neither these files nor the new draft in the available workspace or saved branch tree. Their location and recoverability remain unknown.

Tables~\ref{tab:positive_summary} and \ref{tab:positive_paid} therefore reproduce historical reports, not results independently reverified from currently available raw files. No missing row, certificate or source is reconstructed from the aggregate. The current mathematical specification is available, but equivalence to the earlier executed code cannot be checked. This prevents a complete reproducibility or submission-ready claim.

\subsection{Reported comparison with the complete route}

The checkpoint describes fixed binary $A/B/K$, outward logarithms of positive observations, posterior residual enclosures, MPFR 4.2.1 at 192 bits and a separate reader. It reports both methods accepted on all 43 states in three rotating-order repetitions: 258 dependent method runs. The direct route always constructs complete projections; the local route first uses the stability bound. Both use the same selection tolerance.

\begin{table}[htbp]\centering\small
\caption{Historical aggregate; original execution files are unavailable. Times sum casewise medians across three repetitions and include the separate reader.}\label{tab:positive_summary}
\begin{tabular}{lrr}\toprule
Quantity & Direct route & Local route\\\midrule
Accepted states per repetition & 43/43 & 43/43\\
Complete projections, first repetition & 114 & 25\\
Conditional endpoint evaluations & 1724 & 1372\\
Integration refinements & 544 & 646\\
Checked numerical stage (s) & 10.529803 & 10.229626\\
\bottomrule\end{tabular}\end{table}

The reported projection count falls by 78.1\%, but integration refinements increase. The checked-stage time ratio is only 1.029. A separately recorded posterior preparation took 15.76 s once; I/O, historical coefficient formation and legacy-domain acquisition are excluded from the table. This is a modest dependent-development observation, with no significance, universal speedup, cold-start or physical deployment claim. Earlier solver-only ratios have weaker accounting scopes and do not replace this result.

The checkpoint also reports all 228 nominal branch/method reference values inside the computed intervals and no certificate override in these engineering states. These remain historical reports pending recovery of their underlying files. Repeated numerical calls do not add engineering samples.

\subsection{Reported paid engineering replay}

Decisions were reported frozen before this calculation's paid-trajectory replay. Both checked methods reportedly matched them, allowing reuse of the already paid first and next responses under the original mean continuation. This is dependent descriptive replay, not a new closed-loop experiment. A tie is within $10^{-6}m_0$.

\begin{table}[htbp]\centering\small
\caption{Historical two-evaluation outcome aggregate; original positive-scope decisions and replay files are unavailable. Positive normalized compliance change is worse.}\label{tab:positive_paid}
\begin{tabular}{lrrrr}\toprule
Comparator & Better & Tie & Worse & Mean change/$m_0$\\\midrule
Raw-response mean & 1 & 41 & 1 & $+0.0074365$\\
Raw-response EI & 2 & 40 & 1 & $-0.0045213$\\
Positive-model one-step EI & 1 & 40 & 2 & $+0.0115273$\\
\bottomrule\end{tabular}\end{table}

All median changes are zero. These mixed outcomes do not establish engineering-policy superiority. Reliable solution of the declared model-based selection problem and improvement of the physical objective are distinct empirical questions.

\subsection{Adverse results and fixed resource boundary}

The earlier raw-Gaussian direct calculation remains separate, with available source and recorded outputs. Negative terminal predictions account for 99.40--99.9997\% of branch values; removing that component changes rankings in 38 of 43 states. Its paid mean comparison gives 1 better, 39 tied and 3 worse; its EI comparison gives 0, 41 and 2. The positive planning model is a later development choice prompted by this support mismatch, not a held-out confirmation.

The archived action-set study records identical nominal and band indices in all 228 branch/domain cases, although 908 intervals widen. Its coefficient/filter/check components give a recorded 1.36-fold integer-envelope improvement over the optimized pairwise control. This concerns a different output and cannot be combined with the rollout ratio. The H1 obstruction, failed independent gate and unimplemented H5 extension remain unchanged.

The two permitted witnesses are exhausted. Their true responses, approximately 14324.74 and 9639.92, greatly exceed trial bounds of approximately 232.53 and 200.60. The weak-bound diagnosis is retained without repeating the screen. The project ledger remains 55,309 physical attempts. The historical positive scope records 86 root-system constructions and 774 dependent numerical method runs across three accounting contracts, with zero new physical evaluations, bound queries, hyperparameter fits, retries, replacements or independent optimization samples. This manuscript integration performs none of those scientific runs.
